\subsection{Historical Pipeline and Shared-Proposal Comparison}
\label{sec:eval:ablation}

Table~\ref{tab:repair_ablation} compares the stored component settings.
Without LLM completion, the output is the preprocessed graph. Direct LLM
uses its own prompt and skips preprocessing. The gate comparison scores the
same proposal before and after filtering. Direct LLM is a pipeline comparison,
not an isolated preprocessing switch. Section~\ref{sec:eval:mechanisms} separates
prompt context from the other pipeline choices;
Section~\ref{sec:eval:individual_ablation} tests individual filter and optimizer components.

\begin{table}[t]
\centering
\caption{Controlled pipeline settings. The Direct LLM row also changes the prompt.}
\label{tab:repair_ablation}
\small
\begin{tabular}{l c c c c}
\toprule
Variant & Repair & Over-repair & F1 & Exact \\
\midrule
No LLM completion & 0.3200 & 0.0000 & 0.8472 & 0.0622 \\
Direct LLM (separate prompt) & 0.9244 & 0.0758 & 0.9520 & 0.8267 \\
No constraint gate & 0.9800 & 0.0227 & 0.9895 & 0.9333 \\
\textbf{Diagnosis + Gate} & \textbf{0.9800} & \textbf{0.0167} & \textbf{0.9917} & \textbf{0.9333} \\
\bottomrule
\end{tabular}
\end{table}

LLM completion accounts for most recovered fields. Diagnosis + Gate reaches
99.17\% F1, compared with 95.20\% for Direct LLM. On the shared proposals,
filtering leaves defect repair unchanged and reduces over-repair by
0.60 points. It removes eight unsupported triples, all absent from the
reference graphs.
